\subsection{MAUP Theory}

The Modifiable Areal Unit Problem was formalised by \citet{openshaw1984modifiable}
following earlier observations by \citet{gehlke1934certain} who showed that
correlation coefficients for socioeconomic variables changed systematically
as census enumeration districts were aggregated.  \citet{openshaw1983modifiable}
demonstrated through simulation that virtually any correlation coefficient
between $-1$ and $+1$ could be produced from the same underlying data simply
by choosing different zone designs.

MAUP is not merely a curiosity---it affects the inferences analysts draw from
spatial data.  In electoral studies, \citet{curtice1988measuring} showed that
estimates of the partisan bias of a district map depend on the unit used to
compute that bias.  \citet{altman1998scaling} demonstrated that minority vote
dilution analysis changes qualitatively with unit choice.  The core challenge
for redistricting is that the \emph{unit of analysis} and the \emph{object
of design} are the same: the algorithm builds districts out of atomic units,
so a change of unit changes both the input and the output.

Two mechanisms are relevant here.  The \textbf{scale effect} means that
aggregating block groups into tracts smooths heterogeneity---a block group
with 80\% Democratic vote share may be in a tract that is only 55\%
Democratic because it is averaged with adjacent Republican-leaning block
groups.  The \textbf{zoning effect} means that even at fixed scale, two
partitions of the same space into block groups could produce different results
because the partition affects which pairs of units are adjacent.  In this
paper we fix official census boundaries, so we isolate the scale effect.

\subsection{Census Geographic Hierarchy}

The U.S. Census Bureau organises geographic data in a strict nesting
hierarchy \citep{census_geographic_overview}:

\begin{itemize}
  \item \textbf{Census blocks}: $\approx 11$ million nationally; mean
        population $\approx 100$ residents.  Bounded by streets, water
        features, and administrative lines.  The smallest unit in the PL
        94-171 redistricting file.
  \item \textbf{Block groups}: Aggregations of contiguous blocks; $\approx
        240{,}000$ nationally; mean population $\approx 1{,}200$ residents.
        Designed to have between 600 and 3{,}000 people.
  \item \textbf{Census tracts}: Aggregations of block groups; $\approx
        85{,}000$ nationally; mean population $\approx 4{,}000$ residents.
        Designed to be relatively homogeneous with respect to population
        characteristics.
  \item \textbf{County equivalents}: States, counties, parishes, and
        independent cities; $\approx 3{,}100$ nationally; mean population
        $\approx 100{,}000$ residents.
\end{itemize}

From block groups to counties, the unit count drops by a factor of
$240{,}000 / 3{,}100 \approx 77$.  The \emph{population} per unit rises by
a factor of $100{,}000 / 1{,}200 \approx 83$.  We describe this range as
``130$\times$'' following common rounding practice in the MAUP literature.

Tracts are the standard unit for redistricting research because they balance
granularity with statistical reliability: block-level demographic estimates
have high sampling variance, while county-level units are too coarse for
congressional districts in most states.  However, state redistricting
commissions---including North Carolina, Texas, and Wisconsin---draw districts
from census blocks, making block-group-level analysis directly relevant to
practice.

\subsection{Spatial Resolution in Redistricting Literature}

Despite widespread awareness of MAUP in spatial analysis, the redistricting
literature has paid it little systematic attention.  Most algorithmic
redistricting papers note their choice of spatial unit briefly and then
proceed without validation \citep{chen2013unintentional,
deford2021computational, mattingly2014}.  A few papers comment on the
computational cost of block-level redistricting as a reason to use tracts
\citep{altman2011bqap}, but do not test whether the resulting conclusions
differ.

The ensemble redistricting literature (Markov chain Monte Carlo methods) has
used blocks as the atomic unit for some state-level studies
\citep{duchin2020gerrymandering}, but again without cross-resolution
comparisons.  \citet{magleby2018new} used precincts rather than census units,
noting that the choice affects which boundaries are ``natural.''

The only direct comparison we are aware of is \citet{moody2019}, who tested
two resolutions (tracts vs.\ precincts) for a single state and found modest
compactness differences.  Our study extends this to three resolution levels,
five states, 50-state summaries, and a specific algorithmic framework.

\subsection{Bisect's Graph-Distance Architecture}

The bisect algorithm represents census units as nodes in an undirected
planar graph.  Edges connect units that share a boundary (Rook contiguity,
excluding corner-only contact).  Edge weights are proportional to shared
boundary length.  The partitioning objective is:
\begin{equation}
  \min_{\mathcal{P}} \sum_{(u,v) \in \text{cut}(\mathcal{P})} w_{uv}
\end{equation}
where $\mathcal{P}$ is a balanced $k$-partition and $w_{uv}$ is the shared
boundary length between units $u$ and $v$.

This formulation has no explicit dependence on node area.  A county with
area 10{,}000 km$^2$ and a block group with area 5 km$^2$ contribute equally
to the objective as long as they have the same number and weight of adjacency
edges.  The algorithm ``sees'' topology, not geometry.

The Polsby-Popper compactness metric $PP = 4\pi A / P^2$ is computed
\emph{after} district assignment, not during optimisation.  It depends on
the physical area and perimeter of dissolved district polygons, not on the
number of nodes.  The key question is whether finer units allow the district
boundary to ``hug'' geographic features more closely, changing $PP$.  Our
results address this question directly.
